\section{Introduction}
\label{sec:introduction}

Counterfactual regret minimization (CFR) computes approximate equilibria in
large imperfect-information games by alternating counterfactual evaluation
with local regret updates \citep{zinkevich2007cfr}.  Public chance often
creates a finite cut of outcomes observed by both players.  Evaluating every
outcome gives an exact update; evaluating a smaller batch lowers the work per
update and permits more frequent learner transitions.  Public Chance Sampling
and Monte Carlo CFR (MCCFR) obtain guarantees by making each sampled update a
conditionally valid estimator of its target
\citep{lanctot2009mccfr,johanson2012pcs,farina2020stochastic}.

Many high-throughput solvers use a different pattern.  They process a fixed
public order without replacement, update the learner after each batch, and
continue the order across the epoch.  Every outcome appears equally often,
but later batches are evaluated at profiles produced by earlier feedback.
The delivered row is therefore coupled to its own history.  Exact counts at
the end of an epoch do not make the intermediate rows conditionally unbiased,
and regret on those rows need not control exploitability in the full game.

The two-outcome case isolates the issue.  Let $X_A(\sigma)$ and
$X_B(\sigma)$ be centered public components at profile $\sigma$.  Processing
$A$ and then $B$ delivers $2X_A(\sigma_t)$, updates the strategy, and then
delivers $2X_B(\sigma_{t+1})$.  Relative to dense feedback at both origins,
the epoch error is
\[
 [X_A(\sigma_t)-X_A(\sigma_{t+1})]
 -[X_B(\sigma_t)-X_B(\sigma_{t+1})].
\]
Coverage cancels static selection error, but not the interaction between
component identity and learner motion.  We call the resulting one-sided
full-game correction the \emph{target debit};
Figure~\ref{fig:two-outcome-mechanism} illustrates the mechanism.

\begin{figure}[t]
\centering
\input{figures/two_outcome_mechanism}
\caption{Balanced selection cancels static component error.  Learner updates
between partial batches leave target debit; complete coverage evaluates all
components at one origin and has zero debit.}
\label{fig:two-outcome-mechanism}
\end{figure}

Our central result turns this error into a decision account.  For each local
row, only cumulative coordinates on which delivered feedback
\emph{underestimates} the dense target can increase target regret.
Theorem~\ref{thm:weighted-target-debit} combines this one-sided row transfer
with the native alternating CFR reduction.  It bounds weighted exploitability
by delivered-row regret, target debit, and a native alternating correction.

Debit-Budgeted CFR (DB-CFR) uses this account causally.  At an aligned block
boundary, it selects a width and order from completed history, executes that
plan once, and settles its charge afterward.  Partial execution continues
while the settled account is solvent and at most $\ell$ audits remain
unsettled; otherwise the method executes complete coverage, whose debit is
exactly zero.  Theorem~\ref{thm:debit-controller} shows that hidden debit is at
most the sum of the $\ell$ largest predictable block envelopes.  For fixed
blocks and polynomial output weights, a fixed audit cap contributes only
$O(T^{-1})$ and preserves the additive-RM/RM+ rate.  Thus target validity does
not require either per-event unbiasedness or future candidate replay.

The settled charge must contain semantic information.  Our finite converse compares
two four-round plans with the same public counts, work, task graph, prefix
discrepancy, and scalar strategy-path length.  Their exact controller charges
are $0$ and $11/8$.  Any sound rule based only on those shared summaries must
charge both by at least $11/8$ and therefore rejects the safe plan at budget
one.  A strict-budget admission rule limited to the same summaries must also
reject both before execution.  DB-CFR resolves this information boundary by
settling semantic debit after the selected block and reserving for the bounded
outstanding account.

The empirical program first establishes the decision geometry, then changes
the learner path.  Released HUNL turn endgames show that partial width dominates
at shallow budgets, that order changes quality at fixed coverage, and that
complete coverage overtakes every partial width between 32 and 64 full-cut
budgets.  In public Leduc, widths 2, 5, and 10 occupy one work--debit frontier.
A lagged completed-history rule fixed on four development orders then beats the
best fixed width at no greater installed recurrence work on all ten unseen
orders, with median NashConv improvement 0.0456 and no future candidate
replay.  A separate HUNL resource audit measures the exact selected-path
certificate rather than treating it as free work.

Our contributions are threefold.  \textbf{(i)} We establish a weighted,
one-sided target-debit theorem for deterministic partial public feedback under
native alternating CFR.  \textbf{(ii)} We introduce causal DB-CFR and prove
that bounded settlement removes future candidate replay while preserving the
additive-RM/RM+ rate; strict separations distinguish its charge from symmetric
and metadata-only alternatives.  \textbf{(iii)} We provide two-domain
evidence for the public width/order problem, finite trace certificates, and a
held-out dynamic-width consequence.
